\documentclass[10pt]{article}
\usepackage[T1]{fontenc}
\usepackage[utf8]{inputenc}
\usepackage[french]{babel}
\usepackage{amsmath,amssymb,mathrsfs,mathpazo}
\usepackage{geometry}
\geometry{a4paper,margin=2cm}
\usepackage{graphicx,array,tabularx,tabularray,colortbl,multirow,multicol,booktabs,enumitem,pifont,ulem,fancybox,diagbox,cancel}
\usepackage{pstricks,pstricks-add,pst-plot,pst-tree,pst-node,pst-text,pst-func,pst-math,pst-bspline,pst-3dplot,pst-eucl,pst-solides3d}
\let\pscalculate\undefined
\usepackage{pst-bezier}
\usepackage{tikz,pgfplots,tkz-tab}
\usetikzlibrary{arrows,arrows.meta,patterns,calc,positioning,trees,decorations.markings,shapes.geometric,shapes.misc,shapes.arrows,intersections,decorations.pathmorphing,decorations.pathreplacing,backgrounds,shadows,babel}
\pgfplotsset{compat=1.18}
\usepgfplotslibrary{fillbetween}
\usepackage[np]{numprint}
\usepackage{listings,makecell,pas-tableur,twoopt}
\let\barre\undefined
\usepackage[pstricks]{tablvar}
\definecolor{darkgreen}{RGB}{0,100,0}
\providecommand{\red}{\color{red}}
\providecommand{\blue}{\color{blue}}
\providecommand{\green}{\color{green}}
\providecommand{\black}{\color{black}}
\providecommand{\white}{\color{white}}
\providecommand{\gray}{\color{gray}}
\providecommand{\cyan}{\color{cyan}}
\providecommand{\magenta}{\color{magenta}}
\providecommand{\yellow}{\color{yellow}}
\providecommand{\violet}{\color{violet}}
\providecommand{\orange}{\color{orange}}
\providecommand{\euro}{\texteuro}
\providecommand{\EUR}{\texteuro}
\providecommand{\ieme}{\textsuperscript{e}}
\providecommand{\degres}{\ensuremath{^\circ}}
\providecommand{\degree}{\ensuremath{^\circ}}
\providecommand{\vv}[1]{\overrightarrow{#1}}
\providecommand{\Coord}[1]{#1}
\providecommand{\decosix}{}
\providecommand{\starredbullet}{\ensuremath{\bullet}}
\pagestyle{empty}
\setlength{\parindent}{0pt}
\usepackage{adjustbox}\DeclareUnicodeCharacter{2236}{:}\DeclareUnicodeCharacter{2190}{\ensuremath{\leftarrow}}
\providecommand{\R}{\mathbb{R}}
\providecommand{\N}{\mathbb{N}}
\providecommand{\Z}{\mathbb{Z}}
\providecommand{\Q}{\mathbb{Q}}
\providecommand{\D}{\mathbb{D}}
\providecommand{\C}{\mathbb{C}}
\providecommand{\e}{{\mathrm{e}}}
\providecommand{\vect}[1]{\overrightarrow{#1}}
\providecommand{\vectt}[1]{\overrightarrow{\mathrm{#1}}}
\providecommand{\vvt}[1]{\overrightarrow{\mathrm{#1}}}
\providecommand{\np}[2][]{{#2~#1}}
\providecommand{\vv}[1]{{\overrightarrow{#1}}}
\renewcommand{\d}{{\mathrm{d}}}
\providecommand{\degres}{{^\circ}}
\providecommand{\degre}{{^\circ}}
\providecommand{\blue}{\color{blue}}
\providecommand{\red}{\color{red}}
\providecommand{\enskip}{\quad}
\providecommand{\llbracket}{\lbrack\!\lbrack}
\providecommand{\rrbracket}{\rbrack\!\rbrack}
\providecommand{\euro}{{\text{€}}}
\providecommand{\no}{{\text{n°}}}
\providecommand{\interval}[2]{{[#1;#2]}}
\providecommand{\intervFO}[2]{{[#1;#2[}}
\providecommand{\dotfill}{\ldots}

\providecommand{\ds}{\displaystyle}
\providecommand{\textding}[1]{#1}
\providecommand{\Oij}{\(\mathrm{O};\vec{\imath},\vec{\jmath}\)}
\providecommand{\Oijk}{\(\mathrm{O};\vec{\imath},\vec{\jmath},\vec{k}\)}
\providecommand{\Ouv}{\(\mathrm{O};\vec{u},\vec{v}\)}
\providecommand{\intervFF}[2]{\ensuremath{[#1;#2]}}
\renewcommand{\euro}{\ifmmode\text{\texteuro}\else\texteuro\fi}
\DeclareUnicodeCharacter{00B0}{\ifmmode^{\circ}\else\textdegree\fi}
\DeclareUnicodeCharacter{20AC}{\ifmmode\text{\texteuro}\else\texteuro\fi}
% Énoncés : Éducation nationale. Transcription et corrigés : APMEP, Denis Vergès et contributeurs. Découpage et classement : CoopMaths, Sébastien Lozano, Éric Elter et contributeurs. Adaptation de la présentation : Maths Explorer.
% CC BY-SA : https://coopmaths.fr/www/about/
% https://coopmaths.fr/alea/static/e3c/2020/tex/e3c_2020_00_sujet36_4_cor.tex
\begin{document}
\noindent\textit{APMEP / CoopMaths. CC BY-SA. Présentation adaptée par Maths Explorer.}\par\medskip
\(f(x) = 2x^3 + 2x^2 - 2x + 1\).

\subsection*{1.}

\(f\) fonction polynôme est dérivable sur \(\mathbb{R}\), donc sur \([-3\,;\,3]\) et sur cet intervalle :
\[
f'(x) = 6x^2 + 4x - 2 = 2(3x^2 + 2x - 1)
\]

\subsection*{2.}

\(f'(x)\) a donc le signe du trinôme \(3x^2 + 2x - 1\). Celui-ci a une racine évidente \(-1\) et comme le produit des racines est égal à \(\dfrac{c}{a} = -\dfrac{1}{3}\), l'autre racine est égale à \(\dfrac{1}{3}\).

On sait que ce trinôme est positif sauf sur l'intervalle \(\left[-1\,;\,\dfrac{1}{3}\right[\).

\subsection*{3.}

De la question précédente on déduit que, sur l'intervalle \([-3\,;\,3]\) :
\[f(-3) = 2 \times (-3)^3 + 2 \times (-3)^2 - 2 \times (-3) + 1 = -29 \quad ; \quad f(-1) = -2 + 2 + 2 + 1 = 3\]
\[f\left(\dfrac{1}{3}\right) = 2 \times \dfrac{1}{27} + 2 \times \dfrac{1}{9} - 2 \times \dfrac{1}{3} + 1 = \dfrac{17}{27} \approx 0{,}63 \quad ; \quad f(3) = 54 + 18 - 6 + 1 = 67\]
\begin{center}
\begin{tikzpicture}
\tkzTabInit[lgt=2, espcl=2.5]{$x$ / 1, {$f$} / 2}{${-3}$, ${-1}$, ${\frac13}$, ${3}$}
\tkzTabVar{-/{$-29$},+/{$3$},-/{$\frac{17}{27}\approx 0{,}63$},+/{$67$}}{/}
\end{tikzpicture}
\end{center}

\subsection*{4.}

\paragraph{a.} On sait que :
\[
M(x\,;\,y) \in \mathcal{T} \iff y - f(0) = f'(0)(x - 0).
\]
Avec \(f(0) = 1\) et \(f'(0) = -2\), on obtient :
\begin{align*}
&M(x\,;\,y) \in \mathcal{T} \\
\iff &y - 1 = -2x \\
\iff &y = -2x + 1.
\end{align*}

\paragraph{b.} Un point de \(\mathcal{T}\) appartient à la courbe \(\mathcal{C}\) si ses coordonnées \((x\,;\,y)\) vérifient les deux équations suivantes :
\[
\begin{cases}
y = -2x + 1 \\
y = 2x^3 + 2x^2 - 2x + 1
\end{cases}
\]
Cela donne :
\begin{align*}
&-2x + 1 = 2x^3 + 2x^2 - 2x + 1 \\
\iff &2x^3 + 2x^2 = 0 \\
\iff &2x^2(x + 1) = 0,
\end{align*}
d'où :
\[
\left\{
\begin{array}{l}
2x^2 = 0 \\
x + 1 = 0
\end{array}
\right.
\iff
\left\{
\begin{array}{l}
x = 0 \\
x = -1
\end{array}
\right.
\]
On retrouve donc le point \(A\) d'abscisse \(0\) et le point \(B(-1\,;\,3)\).

\end{document}
